% !TEX root = ../ctufit-thesis.tex
%---------------------------------------------------------------
\chapter{Reproducibility}
%---------------------------------------------------------------

This appendix lists the commands used to regenerate the CSV data, charts, tables,
and thesis figures referenced in the thesis. Commands are intended to be run from
the repository root after the Python environment has been prepared and the
exported mobile and PC benchmark runs are present under
\texttt{benchmarks/results}.

\section{Combined CSV, chart, and table generation}

The main reproducibility command regenerates the combined CSV, chart suite, and
summary tables used by the results chapter:

\begin{verbatim}
./benchmarks/analysis/run_full_analysis.sh \
  --mobile-root benchmarks/results/mobile \
  --pc-root benchmarks/results/pc/runs \
  --combined-csv benchmarks/results/combined/combined_runs_from_json.csv \
  --charts-dir benchmarks/results/combined/charts \
  --tables-dir benchmarks/results/combined/tables \
  --focus-device "xiaomi redmi note 8 pro" \
  --devices "iphone 14,xiaomi redmi note 8 pro" \
  --max-resolved-batch 8
\end{verbatim}

If the AFW landmark annotation directory is available outside the attachment,
pass it explicitly so mobile landmark accuracy can be recomputed from exported
predictions:

\begin{verbatim}
./benchmarks/analysis/run_full_analysis.sh \
  --afw-dir path/to/afw \
  --focus-device "xiaomi redmi note 8 pro" \
  --devices "iphone 14,xiaomi redmi note 8 pro" \
  --max-resolved-batch 8
\end{verbatim}

\section{Readable thesis chart regeneration}

The thesis uses enlarged, thesis-readable versions of selected charts. After the
combined CSV and summary tables have been regenerated, run:

\begin{verbatim}
./.venv/bin/python \
  benchmarks/analysis/generate_readable_thesis_charts.py
\end{verbatim}

\section{Introductory pipeline figure regeneration}

The introductory face-processing figure is generated from the sample AFW image
and its landmark annotation:

\begin{verbatim}
./.venv/bin/python \
  benchmarks/analysis/generate_thesis_pipeline_figure.py
\end{verbatim}

\section{INT8 XNNPACK landmark model export with post-training quantization}

\begin{verbatim}
EXPORT_BATCH_SIZES=1 \
EXPORT_DYNAMIC_BATCH=1 \
EXPORT_DYNAMIC_MIN_BATCH=1 \
EXPORT_DYNAMIC_MAX_BATCH=128 \
XNNPACK_INT8_CALIBRATION_SAMPLES=32 \
XNNPACK_INT8_METHODS=pc,pt \
./models/export_landmark_models_xnnpack_quantized.sh
\end{verbatim}

Optional MobileFaceNet stability switch (enabled by default):

\begin{verbatim}
MOBILEFACENET_STATIC_INT8_DYNAMIC_ACT_WORKAROUND=1
\end{verbatim}
